\begin{foldable} % LLM --> DD
  The progress of large language models depends on large-scale supervised corpora, yet the cost of curating and training on such data compounds across fine-tuning cycles, hyperparameter sweeps, and continual learning.
  The real burden is not any single run but the cumulative expenditure across the full development lifecycle.
  Corpus size, not model size, is often the binding constraint.
  Reducing it is therefore a systemic priority with direct implications for cost, carbon footprint, and data governance.
\end{foldable}

\begin{foldable} % DD --> Text DD
  Dataset distillation (DD), first proposed by~\cite{wang2018dataset}, replaces a large corpus $\mathcal{D}$ with a smaller surrogate $\tilde{\mathcal{D}} \ll \mathcal{D}$ such that a model trained on $\tilde{\mathcal{D}}$ performs comparably to one trained on $\mathcal{D}$.
  Put another way: what is the smallest textbook that teaches the same exam?
  This places DD squarely within data-centric AI, complementing model compression by operating on the data rather than the model itself.
  Early DD methods work in image space, optimizing synthetic pixels via gradient or trajectory matching~\cite{zhao2020dataset,cazenavette2022dataset}.
  Text breaks these assumptions: non-differentiable token decoding blocks gradient-based synthesis, and embedding spaces are tightly coupled to specific architectures.
  Adapting DD to text has therefore required fundamentally different strategies.
\end{foldable}

\begin{foldable} % Baseline text DD methods -- gap --> desiderata --> method
  Existing text DD methods~\cite{sucholutsky2021soft,li2021data,maekawa2023dataset,maekawa2024dilm,tao2024textual} fall short in at least one of the following: at extreme compression, uniform weighting exhausts the budget on uninformative samples; their objectives lack task alignment or global optimization; and their outputs are embeddings, neither auditable nor transferable.
  We survey this landscape in \S\ref{sec:02-related-textdd} and show that no prior method simultaneously addresses all of the above.
  These gaps motivate three desiderata for TAKE:{
    (1)~\emph{knowledge-based weighting} --- weighting samples by their task-aligned downstream contribution, directing the budget toward the most informative ones;
    (2)~a \emph{dataset-level distillation objective} --- optimizing over the full training distribution rather than batch-by-batch, ensuring global coverage;
    and (3)~\emph{human-readable output} --- a distilled corpus that is directly inspectable, auditable, and transferable across architectures.
  }
\end{foldable}

\begin{foldable} % Our method TAKE and contributions
  We propose \textbf{TAKE} (Trajectory-Aware Knowledge Estimation), the first text DD method to satisfy all three desiderata.
  We address the first via influence functions~\cite{koh2017understanding}, which quantify each sample's counterfactual effect on the downstream loss and provide a theoretically grounded, task-aligned weighting scheme (\S\ref{sec:03-theory}).
  However, scoring at a single checkpoint introduces the \textit{hard-sample bias} (HSB): at convergence, noisy samples dominate influence scores while clean and moderate samples --- the most valuable for a robust distilled dataset --- are suppressed.
  TAKE corrects this by integrating influence scores along the full training trajectory into a per-sample knowledge score that captures clean and moderate samples while down-weighting noisy ones.
  Then, TAKE fine-tunes an LLM to generate a pool of human-readable candidate instances.
  The knowledge scores and synthetic candidates jointly feed a discrete Optimal Transport (OT) objective that globally selects the distilled corpus.
  Concretely, we contribute:
  \begin{itemize}
    \item {
        \textbf{Theory:} We formalize the hard-sample bias~(HSB) in the text DD setting, prove that single-checkpoint influence scores are biased toward noisy samples, and derive a formal gap bound for knowledge-reweighted distribution matching.
      }
    \item {
        \textbf{Method:} TAKE is the first text DD method to jointly address sample weighting and distributional coverage --- via trajectory-integrated knowledge scores and a discrete OT objective aligned with the DD goal.
      }
    \item {
        \textbf{Empirical:} TAKE matches or exceeds prior text DD methods across five language tasks at extreme compression (0.1\% or 20 instances/class), producing human-readable distilled corpora that generalize across diverse backbone families.
      }
  \end{itemize}
\end{foldable}
